\section{Corpus Provenance and Claim-Boundary Narrative}
\label{sec:oc13-publication-boundary}

This chapter fixes the public claim boundary of Core~1.3.3. Its task is to state
which scientific authority is canonical, how source witnessing is enforced, how
theorem fate is separated across outward artifacts, and how later channels must
remain aligned with the source corpus, science-state register, and flagship monograph. It does not re-argue the formal doctrine or
theorem spine; those burdens belong to the surrounding chapters and appendices.

\subsection{Canonical scientific authority}

The public scientific authority of the release is this monograph together with the curated evidence package and checksum-bound evidence records. The source-owned science corpus at
evidence-package record science sources is the
authoring and witness locus from which that science-state register is projected. Historical
source witnesses, including earlier Core~1.2/Core~1.3.3 repair baselines, remain
provenance anchors; they are not the current projection authority. The present
master monograph is the flagship reader-facing projection and compiled
reference volume derived from those records. It keeps the orientation ladder,
formal doctrine, theorem route, empirical discipline, synthesis, practical
utility, and audit layers visible together, but it does not displace the science-state register
as the final authority on claim scope, promotion state, or proof boundary.

This rule matters because foundational programmes accumulate several different
textual layers over time: source-native definitions and theorems, explanatory
prose, repair notes, channel-specific summaries, and later operational packets.
Core~1.3.3 treats those layers as distinguishable, not interchangeable.

\subsection{Source witnessing}

Core~1.3.3 is source-first in the strict sense that the flagship claims used in
public release must remain bound to public source witnesses rather than to late
summary prose alone. The point is not archival fetishism. The point is to stop
drift between five things that often blur inside large research programmes:

\begin{itemize}
\item theorem-bearing statements,
\item structural and interpretive explanation,
\item repair and release notes,
\item condensed public packets,
\item and later operational projections.
\end{itemize}

When that boundary blurs, a summary sentence can begin to masquerade as the
theorem itself, or a release dashboard can begin to replace the actual
scientific record. Core~1.3.3 blocks that failure by keeping source witnessing
bound to the source-owned science corpus and the internal science-state register projection
rather than relegating it to hidden control files.

\subsection{Theorem fate and scope discipline}

The present release also separates theorem fate explicitly. Not every inherited
row serves the same role in every outward artifact.

\begin{itemize}
\item bounded theorem route rows may carry positive proof-routed claims
      when their proof and dependency route remains exact.
\item bridge-only support row rows may remain scientifically useful under
      explicit scope notes, but they cannot be promoted as proof-routed.
\item support-only row and future-work support row rows belong in support
      layers, appendices, narrower companion artifacts, or future work rather
      than in the positive body of every outward paper.
\item negative-control rejection row rows are excluded from the positive proof path. They may
      be mentioned only as negative audit history or as a falsifier record, not
      as live support.
\item Frontier or hypothesis material may be discussed only under visibly
      labelled frontier language and may not be silently promoted by summary
      prose.
\end{itemize}

The public positive closure claim is scoped exactly to
current bounded science surface. support-only row rows may still close as
frame-level bounded evidence recorded rows inside the synthesis surface, but they do not become
standalone bounded theorem-routed replay promotions. future-work support row and
negative-control rejection row rows remain outside the public bounded evidence recorded scope for different
reasons: future-work support row rows are frontier/support lanes awaiting later
promotion or rejection, while negative-control rejection row rows are negative-audit lanes.
Neither class can carry positive closure claims in this release.
For avoidance of doubt, a support-only row bounded evidence recorded row is excluded
from the public positive closure claim unless a later release promotes it into
current bounded science surface with the required proof-routed support.

This is not retreat. It is a condition of honest publication. The monograph
can remain large without rhetorical inflation only if the reader can tell which
rows are theorem-bearing, which are bounded, and which are support or frontier
material.

\subsection{Relation to journal-core and companion artifacts}

Core~1.3.3 therefore adopts a one-way reader-facing publication logic:
\[
\begin{gathered}
\text{Source-owned science corpus}\\
\text{plus internal science-state register}\\
\Downarrow\\
\text{Master Monograph}\\
\Downarrow\\
\text{Derived outward artifacts}
\end{gathered}
\]
This diagram describes publication order, not the full provenance graph of
the release package. Several editorial JSON surfaces, generated science
chapters, and release matrices are direct science-state register projections or mirrors of other
governing canonical surfaces. They do not become true merely because the
master monograph cites them, and they are not descendants of the monograph
alone. The monograph is the flagship reading object; the source-owned science
corpus remains the authoring and witness locus, and the science-state register remains the canonical
projection authority for release surfaces.

The journal-core article is narrower by design. Companion dossiers, matrices,
and release surfaces are narrower or more operational by design. Those
artifacts may compress, subset, or reorganise the flagship material for a
specific task, but they may not enlarge the scientific claim, widen notation,
or outrun the theorem-fate discipline fixed here.

The practical consequence is simple. If an outward reader-facing artifact says
something stronger than the source corpus and science-state register authorise, the outward
artifact is wrong, even if it is shorter, cleaner, or easier to market. The
flagship boundary is the publication rule that enforces that ceiling; it is
not an additional source of scientific authority.

\subsection{Downstream projection rule}

The same boundary stabilizes later use in review, collaboration, product,
business, and public communication. Different channels may legitimately cite
different parts of the flagship volume:

\begin{itemize}
\item the theorem route when bounded proof is the issue,
\item the empirical route when promotion discipline is the issue,
\item the practical route when use and comparison are the issue,
\item the audit appendices when challenge and traceability are the issue.
\end{itemize}

What they may not do is invent a second scientific authority. Source witnessing,
theorem fate, and public claim boundary exist so that external use can remain
concise without becoming interpretively lawless.

\subsection{External criticism closure route}
\label{sec:oc13-source-audit-external-criticism-closure}

The Stanislav Tsukrov external review package is closed through Appendix~the oc13 external criticism closure discussion. The route records which objections were answered by source-tradition repair, proof/evidence rows, replay checks, theorem-fate discipline, tuple/K rigor evidence, authorization-bounded journal submission requires separate approval boundaries, and duplicate mappings. It does not move raw review text into the public release and it does not promote a canonical claim.

\subsection{Source-tradition remediation after external review}
\label{sec:oc13-source-tradition-remediation-v3}

The Stanislav Tsukrov review identified a real scholarly weakness: Core~1.3.3
already cited physics-adjacent complexity literature, but it did not adequately
name the systems, autopoiesis, identity, complexity-measure, tipping-point, and
formal-geometry traditions in which several OC constructs must be judged. This
release therefore lowers the novelty ceiling and makes the inherited ground
explicit.

The hierarchy and systems-language route is now read against general systems
theory and mathematical systems theory, including von Bertalanffy,
Boulding, Miller, Mesarovic--Takahara, and Klir
\cite{bertalanffy_general_system_theory,boulding_general_systems_theory,miller_living_systems,mesarovic_takahara_general_systems,klir_architecture_systems}.
The rebirth and post-collapse identity route is read against autopoiesis and
identity/persistence debates, including Maturana--Varela, Varela, Parfit, and
Wimsatt \cite{maturana_varela_autopoiesis,varela_biological_autonomy,parfit_reasons_persons,wimsatt_reductionism_levels}.
The complexity and early-warning route is read against effective complexity,
predictive information, integrated information, and critical-transition
literature \cite{gell_mann_effective_complexity,bialek_predictability_complexity,tononi_information_integration,scheffer_critical_transitions,dakos_early_warning}.
Formal differential-geometric machinery is treated as inherited mathematical
tooling rather than as an OC novelty claim; Olver is the current compact source
witness for the jet/Lie-group side of that boundary \cite{olver_lie_groups}.

The repaired claim is narrow: OC contributes a typed cross-level interface,
collapse/residue/rebirth grammar, and release discipline. It does not claim to
replace GST, autopoiesis, complexity theory, identity metaphysics, tipping-point
analysis, or formal differential geometry.

\paragraph{Bridge to the next chapter.}
With the public claim boundary fixed, the next task is internal consistency:
which statement classes the flagship volume uses, why the hierarchy runs
through \(K_{12}\), and how upper levels remain lawful without rhetorical
inflation.

\clearpage
